% Part: first-order-logic
% Chapter: completeness
% Section: downward-ls

\documentclass[../../../include/open-logic-section]{subfiles}

\begin{document}

\olfileid{fol}{com}{dls}
\olsection{লোয়েনহাইম--স্কোলেম উপপাদ্য}

লোয়েনহাইম--স্কোলেম উপপাদ্য বলে, কোনো তত্ত্বের একটি অসীম মডেল থাকলে তার
এমন একটি মডেলও আছে যা সর্বাধিক !!{denumerable}। এই তথ্যের একটি তাৎক্ষণিক
পরিণাম হলো, প্রথম-ক্রমের যুক্তিবিদ্যা কোনো !!a{structure}-এর আকার
!!{nonenumerable}—এ কথা প্রকাশ করতে পারে না: সব !!{nonenumerable}
!!{structure}s-তে পরিতৃপ্ত যেকোনো !!{sentence} বা !!{sentence}s-এর সেট
কোনো !!{enumerable} গঠনেও পরিতৃপ্ত হয়।

\begin{thm}
\ollabel{thm:downward-ls} $\Gamma$ সঙ্গত হলে তার একটি !!a{enumerable}
মডেল আছে; অর্থাৎ, এমন কোনো !!a{structure}-এ সেটি পরিতৃপ্তিযোগ্য যার
সংজ্ঞাক্ষেত্র সসীম অথবা !!{denumerable}।
\end{thm}

\begin{proof}
$\Gamma$ সঙ্গত হলে পূর্ণতা উপপাদ্যের প্রমাণ থেকে পাওয়া
!!{structure}~$\Struct M$-এর সংজ্ঞাক্ষেত্র~$\Domain{M}$ ভাষা~$\Lang L$-এর
পদগুলির সেটের চেয়ে বড় নয়। তাই $\Struct M$ সর্বাধিক !!{denumerable}।
\end{proof}

\begin{thm}
\ollabel{noidentity-ls} $\Gamma$ যদি অভিন্নতাবিহীন প্রথম-ক্রমের
যুক্তিবিদ্যার ভাষায় !!{sentence}s-এর একটি সঙ্গত সেট হয়, তবে তার একটি
!!a{denumerable} মডেল আছে; অর্থাৎ, অসীম ও !!{enumerable} সংজ্ঞাক্ষেত্রের
কোনো !!a{structure}-এ সেটি পরিতৃপ্তিযোগ্য।
\end{thm}

\begin{proof}
$\Gamma$ সঙ্গত এবং তাতে অভিন্নতাযুক্ত কোনো বাক্য না থাকলে, পূর্ণতা
উপপাদ্যের প্রমাণ থেকে পাওয়া !!{structure}~$\Struct M$-এর
সংজ্ঞাক্ষেত্র~$\Domain{M}$ ভাষা~$\Lang L'$-এর পদগুলির সেটের সঙ্গে
অভিন্ন। তাই $\Trm[L']$ !!{denumerable} বলে $\Struct{M}$-ও তাই।
\end{proof}

\begin{ex}[স্কোলেমের কূটাভাস]
জার্মেলো--ফ্রেঙ্কেল সেটতত্ত্ব~$\Log{ZFC}$ এমন একটি অত্যন্ত শক্তিশালী
কাঠামো, যেখানে সেটের আকারসংক্রান্ত তথ্যসহ কার্যত সব গাণিতিক বক্তব্য
প্রকাশ করা যায়। যেমন, $\Log{ZFC}$ প্রমাণ করতে পারে যে বাস্তব সংখ্যার
সেট~$\Real$ !!{nonenumerable}; যেকোনো সেটের ঘাতসেট সেই সেটের চেয়ে বড়—
ক্যান্টরের এই উপপাদ্যও সেটি প্রমাণ করতে পারে; ইত্যাদি। $\Log{ZFC}$ সঙ্গত
হলে তার সব মডেল অসীম; আরও, সেগুলির প্রত্যেকটিতে এমন !!{element}s থাকে
যেগুলিকে তত্ত্বটি !!{nonenumerable} বলে—যেমন, যে উপাদানটি স্বাভাবিক
সংখ্যার ঘাতসেটের অস্তিত্ববিষয়ক $\Log{ZFC}$-র উপপাদ্যকে সত্য করে।
লোয়েনহাইম--স্কোলেম উপপাদ্য অনুসারে $\Log{ZFC}$-র !!{enumerable}
মডেলও আছে—এমন মডেল, যার মধ্যে “!!{nonenumerable}” সেট আছে, কিন্তু মডেলটি
নিজেই !!{enumerable}।
\end{ex}

\end{document}
